\section{Isolation and dominated extensions}
\label{iso:section-isolation}

Throughout this section, $\mathcal C$ is a working scheme in the sense of
Definition~\ref{lev:working-package}. We use the refinements of the eligible
levels required in Section~\ref{calc:section-calculus}. All objects and
extensions in a displayed type are at one eligible level, unless a change of
level is stated. We retain the eligible-subclass hypothesis of
Convention~\ref{lev:eligible-subclass}; every larger auxiliary sufficiently
closed level below is chosen in its unbounded refined subclass.
A type over an object at a smaller eligible level is computed
in a section of that level; Corollary~\ref{calc:mixed-realization} makes this
independent of the section. A tuple is bounded when its cardinality is less
than the level of its base. We write $p^{\mathrm{fr}}_B$ for the free lift
$p^B$ of Theorem~\ref{calc:free-extension} when its role needs emphasis.

We will repeatedly use the following convention concerning bounds. A request
involving a tuple of cardinality $\sigma$ is first enlarged to an infinite
arity at least $\sigma$. The testing, compact-placement, and approximation
bounds of Section~\ref{calc:section-calculus} are then applied at that arity.
There are only finitely many such enlargements in each argument below. A
source level is required to be strictly closed under those bound functions.
Thus, when a tuple is bounded at a smaller level, every test used to determine
its type may be taken below that smaller level, even if the receiving object
is at a much larger eligible level. This does not assert saturation over an
unbounded parameter set.

Our goal is an extension controlled by a given mark, with an embedding
into every ambient object realizing that mark's type.  Relative isolation
will give this control on bounded tuples.  The main existence argument
first treats one bound on tuple length and then obtains all bounded
lengths in a single extension.

\subsection{Relative isolation}

\begin{definition}\label{iso:relative-isolation}
Let $q=\tp(u,c/A)$ and $p=\tp(u/A)$, where the concatenation $(u,c)$ is
bounded. We say that $q$ is \emph{isolated relative to $u$} if there is a
bounded set $Z\subseteq A$ such that, for every $q'=\tp(u',c'/A)$,
\[
 \tp(u'/A)=p\quad\text{and}\quad
 D(u',c',Z)=D(u,c,Z)
 \quad\Longrightarrow\quad q'=q.
\]
Here and below a parameter set in a diagram is supplied with a fixed
enumeration, and the same enumeration is used on both sides. We call $Z$ an
\emph{isolating fragment}. A projection retains the entire tuple $u$ unless
explicitly stated otherwise.
\end{definition}

\begin{lemma}\label{iso:free-isolation}\label{iso:upward-transfer}
The following hold at sufficiently closed eligible levels.
\begin{enumerate}
\item A projection of a relatively isolated type is relatively isolated.
\item Suppose $A\leq B$, $q=\tp(u,c/A)$ is isolated relative to $u$, and an
actual coupling over $B$ satisfies
\[
 \tp(u/B)=\tp(u/A)^{\mathrm{fr}}_B.
\]
Then
\[
 \tp(u,c/B)=q^{\mathrm{fr}}_B.
\]
Moreover this type is isolated relative to $u$, using the same isolating
fragment as $q$.
\end{enumerate}
The second assertion also holds when $A$ is at a smaller eligible level than
$B$, provided $(u,c)$ and the required uniform bounds are below the level of
$A$.
\end{lemma}

\begin{proof}
For the first assertion, write $c=(c_0,c_1)$ and let $Z$ isolate $q$.
Suppose $(u',c'_0)$ has the prescribed marginal on $u$ and the same diagram
as $(u,c_0)$ over $Z$. In a model containing $A,u',c'_0$, bounded saturation
extends this diagram by a tuple $c'_1$ so that
\[
 D(u',c'_0,c'_1,Z)=D(u,c_0,c_1,Z).
\]
Isolation identifies the full type over $A$, and its projection identifies
$\tp(u',c'_0/A)$.

For the second assertion, fix an isolating fragment $Z$. It is enough to
compare the actual joint type over $B$ with $q^{\mathrm{fr}}_B$ on the
uniform testing arity for $(u,c)$. Enlarge that arity to include $Z$ and
choose approximation lists in $A$ for both $q$ and $p=\tp(u/A)$, long enough
to evaluate all the comparisons just specified. Let $t$ be a test tuple
from $B$ at this enlarged arity. Bounded saturation moves $t$ to a tuple
$t'$ in $A$, preserving its diagram over $Z$ and both approximation lists.

Choose a realization $u_0$ of $p$ over $A$. The free evaluation of $p$ gives
\[
 D(u,t,Z)=D(u_0,t',Z).
\]
Saturation over the bounded list $(u_0,t',Z)$ supplies $c_0$ such that
\[
 D(u_0,c_0,t',Z)=D(u,c,t,Z).
\]
In particular $(u_0,c_0)$ agrees with $q$ on $Z$ and has marginal $p$.
Isolation gives $\tp(u_0,c_0/A)=q$. Consequently the actual answer to the
query $t$ is the answer that $q$ gives to $t'$. Because the query was moved
over the approximation list for $q$, that is also the free evaluation of
$q$ on $t$. The testing theorem identifies the types over $B$.

Now let a competing type over $B$ have the same marginal on $u$ and agree
with the displayed type on $Z$. Its restriction to $A$ is $q$. The argument
just given identifies it with $q^{\mathrm{fr}}_B$, proving isolation there.
In the mixed-level case, perform these comparisons at the uniform test
bound below the source level. Mixed realization supplies the required
sections, and the same argument applies to those diagrams.
\end{proof}

\begin{proposition}\label{iso:alternative}
There is a bound function $h$, depending only on the fixed working scheme,
with the following property. Put
$\sigma=|u,c|+\aleph_0$. At sufficiently closed eligible levels, precisely
one of the following alternatives holds for $q=\tp(u,c/A)$:
\begin{enumerate}
\item $q$ is isolated relative to $u$ by a fragment of size at most
$h(\sigma)$;
\item there are an extension $A\leq B$ at the same level and an extension
$q_B$ of $q$ such that
\[
 q_B\ne q^{\mathrm{fr}}_B,
 \qquad
 q_B\mathbin{\upharpoonright}u
   =\tp(u/A)^{\mathrm{fr}}_B.
\]
\end{enumerate}
In particular, every relatively isolated type has an isolating fragment
bounded by $h(\sigma)$, uniformly in the level and in the type.
\end{proposition}

\begin{proof}
Choose an infinite $\chi$ large enough for the testing theorem at arity
$\sigma$, for the approximation lists for $q$ and its marginal $p$, and
for the finitely many joint comparisons made below. These choices are
uniform in $\sigma$. Increasing $\chi$ if necessary, all the supports in
the chosen approximation lists have cardinality at most $\chi$.

Present $A=F(I)$ on a well-order of the level cardinality $\mu$. We may
choose $I$ to have more than $\sigma$ disjoint consecutive blocks, each
of order type and cardinality $\mu$: this is still a well-order of
cardinality $\mu$, and all such presentations are objects at this level.
Let $J_0$ consist of the supports of the approximation lists, together
with an initial segment of order type $\chi^+$ in every candidate block.
Its size has a bound depending only on $\sigma$. Let $Z_0$ be the set of
all term values on $J_0$. Since the label set is fixed, the same is true
of $|Z_0|$.

If $Z_0$ isolates $q$, we have the first alternative. Suppose it does not.
Choose $q^1\ne q^0=q$ with the same marginal $p$ and the same diagram on
$Z_0$. By Lemma~\ref{calc:compact-placement}, give both types placements
over the displayed presentation of $A$, and amalgamate their support
orders over $I$. The combined new support has cardinality at most
$\sigma$. Each new support position either hits the interior of one
candidate block or determines an interior cut of at most one such block.
We can therefore choose a candidate block $J$ whose interior contains
neither a hit nor a cut of either placement.

Set $A_*=F(J_0\cup J)\leq A$. This object has cardinality $\mu$. The
restrictions of $q^0$ and $q^1$ to $A_*$ are equal. To check this, take a
test at the fixed sufficient arity. Move every position of the test
lying in $J$ into the initial $\chi^+$ part of $J$, preserving the order
and equality pattern of the test. This can be done above any positions
of that initial part which need to remain in the pattern: fewer than
$\chi^+$ positions have been specified. Leave the positions outside
$J$ unchanged. No position of either candidate separates the interior
of $J$, so the move preserves the natural diagram with either candidate.
The moved query uses only $J_0$ and therefore has the same answer for
$q^0$ and $q^1$. The testing theorem proves equality over $A_*$.

The approximation lists for $q^0$ and $p$ lie in $A_*$. Their majority
evaluations show that $q^0$ and $p$ over $A$ are free from $A_*$. Also,
an isomorphism $f:A_*\simeq A$ may be chosen to fix those bounded lists,
by bounded diagram homogeneity. The transport assertion for free types
then gives
\[
 f\bigl(q^0\mathbin{\upharpoonright}A_*\bigr)=q.
\]
Transport the inclusion $A_*\leq A$ along $f$, viewing its larger object
as an extension $B$ of the original base $A$. Transport $q^1$ as well.
Its restriction is $q$ because the two restrictions to $A_*$ agree. Its
marginal is the free lift of $p$. Its joint type is not free, since the
unique free lift before transport was $q^0$, distinct from $q^1$.
This is the second alternative.

Choose $h(\sigma)$ to bound the size of $Z_0$ for these uniform choices.
Lemma~\ref{iso:free-isolation} excludes the second alternative when $q$
is isolated, even if an initially given isolating fragment was larger.
This proves both exclusivity and the asserted uniform bound.
\end{proof}

\subsection{Primeness and descent of isolation}

\begin{definition}\label{iso:all-dominated-definition}
Let $X\leq Y$ be same-level objects and let $b$ be a bounded tuple in
$Y$. We say that $Y$ is \emph{all-dominated over $X+b$} if, for every
bounded tuple $c$ from $Y$, the type $\tp(b,c/X)$ is isolated relative
to $b$.
\end{definition}

The quantifier over all bounded tuples is part of this definition. It is
not enough to verify isolation for individual entries, or for the tuples
in one chosen countable exhaustion of $Y$.

\begin{proposition}\label{iso:prime}
Let $Y$ be all-dominated over $X+b$ at a sufficiently closed eligible
level $\mu$.
\begin{enumerate}
\item If an extension $N$ of $X$ contains a tuple $b'$ of the same type
over $X$, there is an arrow $Y\longrightarrow N$ over $X$ sending $b$
to $b'$. The target may be at a larger eligible level. A prescribed
bounded tuple match can be retained whenever its joint type with $b$
over $X$ already agrees.
\item Two all-dominated objects over the same data $X+b$ are isomorphic
over those data. The isomorphism may retain a prescribed bounded tuple
match of the kind just described.
\item In an actual coupling, if a bounded tuple $x$ has product with
$b$ over $X$, then $\tp(x/Y)$ is free from $X$.
\end{enumerate}
\end{proposition}

\begin{proof}
Choose increasing bounded lists $c_n$, $n<\omega$, exhausting $Y$;
include $b$, the prescribed initial data, and exhaustions of every vertex
by smaller strong submodels. Suppose the image of $c_n$ has already been
chosen with the correct full type over $X$. Let $Z_{n+1}\subseteq X$
isolate $\tp(b,c_{n+1}/X)$. The old image and the source list have the
same diagram over $Z_{n+1}$. Saturation in $N$ extends this bounded match
to $(b,c_{n+1},Z_{n+1})$. Relative isolation makes its full type over $X$
correct. Taking the union gives a map on all of $Y$ fixing $X$: whenever
an entry of $X$ occurs in a list, its equality with the corresponding
parameter is part of its full type. The strong-submodel exhaustions and
the map criterion in the working package show that this map is an arrow.
This proves the first assertion, also into a larger target because every
saturation request is bounded at the source level.

For the second assertion, perform the same construction alternately on
the two sides, at each step including the next list in the chosen
exhaustion on the side to be covered. Both objects have isolation for
the enlarged bounded list. The construction therefore continues in both
directions and its union is an onto structural isomorphism.

For the third assertion, use symmetry and the definition of product to
choose a model $W\geq X$ containing $x$ over which $b$ is free from $X$.
Type-identifying amalgamation brings $W$ and the actual coupling into a
common extension, identifying $x$ and $b$ over $X$. For every bounded
$c$ in $Y$, Lemma~\ref{iso:free-isolation} says that $(b,c)/W$ is free
from $X$. Hence $x$ and $(b,c)$ have product over $X$. Symmetry and
projection give the free evaluation of $x$ on every bounded test $c$
from $Y$. The testing theorem identifies $\tp(x/Y)$ with the free lift
of $\tp(x/X)$.
\end{proof}

\begin{lemma}\label{iso:downward-transfer}
Suppose $A\leq B$ are objects at eligible levels $\nu\leq\mu$,
$|u,c|<\nu$, and the actual type
\[
 q_B=\tp(u,c/B)
\]
is the free lift of $q_A=\tp(u,c/A)$. If $q_B$ is isolated relative to
$u$, then $q_A$ is isolated relative to $u$. Here $\nu$ is sufficiently
closed for the uniform bounds at $|u,c|+\aleph_0$.
\end{lemma}

\begin{proof}
By Proposition~\ref{iso:alternative}, choose a large-base isolating
fragment $W\subseteq B$ with
\[
 |W|\leq h(|u,c|+\aleph_0)<\nu.
\]
Let $w$ enumerate $W$. The type $\tp(w/A)$ is a small-level type,
computed by taking a level-$\nu$ section containing $A$ and $w$.
Choose a sufficiently long bounded approximation list $(w_i)_{i<\tau}$
for this type inside $A$, where $\tau<\nu$. Enlarge its testing bound
to include the arities of $w$ and $(u,c)$ and choose $\tau$ larger than
the two exception bounds that will occur below. Put
\[
 Z=\bigcup_{i<\tau}\ran(w_i)\subseteq A.
\]
This is a bounded fragment.

Let $q'_A$ have the same marginal on $u$ as $q_A$ and agree with it over
$Z$. Lift $q'_A$ freely to $B$. The marginal on $u$ then agrees with
the marginal of $q_B$. In either free coupling, symmetry identifies the
joint diagram with $w$ by the majority evaluation of $\tp(w/A)$ against
the candidate tuple. Thus all but the prescribed bounded number of
$i<\tau$ give the diagram of that candidate with $w$. Choose one index
outside the two exception sets. The candidates agree on $w_i$, because
$w_i$ lies in $Z$, so their free lifts agree on $w$. Isolation of $q_B$
identifies the free lifts, and restriction identifies $q'_A$ with $q_A$.
Every approximation and comparison here has size below $\nu$; no
approximation to the whole of $B$ is used.
\end{proof}

\subsection{Existence with a prescribed bound on tuple length}

We first isolate the continuity requirement on the marked marginal. This
will let us identify exactly where the proof for a finite mark must be
supplemented for a long independent list.

\begin{definition}\label{iso:marginal-continuity}
A bounded marginal type $p$ satisfies the \emph{continuity condition used
here} if the following holds at every sufficiently closed eligible level
to which $p$ is freely lifted. Suppose
\[
 X_0\leq X_1\leq\cdots\leq X_i\leq\cdots\quad(i<\delta)
\]
is a continuous chain at that level, its union is an object at that level,
and a fixed realization $b$ in a coupled extension satisfies
\[
 \tp(b/X_i)=\tp(b/X_0)^{\mathrm{fr}}_{X_i}
 \qquad(i<\delta).
\]
Then the same equality holds over $\bigcup_{i<\delta}X_i$.
The type at $X_0$ is required to be the appropriate free lift of $p$.
\end{definition}

Every finite-tuple marginal has this property. At a limit, local character
for the finite tuple gives freeness from some earlier stage, and
transitivity gives freeness from $X_0$. This is a direct application of
Proposition~\ref{calc:local-character}, including when the chain has small
cofinality.

\begin{proposition}\label{iso:finite-existence}
Let $\lambda$ be a sufficiently closed eligible level, let $X$ be an
object there, and let $p=\tp(b/X)$ be a finite-tuple type. For every
infinite $\theta<\lambda$ there is a same-level extension $X\leq U$
containing a realization of $p$ such that
\[
 \tp(b,c/X)\text{ is isolated relative to }b
 \quad\text{for every }c\in U\text{ with }|c|\leq\theta.
\]
More generally the assertion holds for a bounded mark $b$ satisfying
Definition~\ref{iso:marginal-continuity}, with $|b|\leq\theta$.
The closure requirements at the output level are uniform functions of
bounded arities; they do not require $\theta$ to have been fixed before
the level was chosen.
\end{proposition}

\begin{proof}
We prove the more general statement. Fix a bounded approximation list
$I_p$ in $X$ which determines $p$ at a sufficient testing arity. Choose
a much larger eligible level $\kappa$ and a cardinal $\rho$ such that
\[
 \lambda<\rho=\rho^\lambda,
 \qquad \rho^+<\kappa.
\]
Require $\rho$ to dominate the approximation and small-hull costs at
arity $\lambda$, and choose $\kappa$ after those set-sized bounds. Such
choices exist: first choose a sufficiently large power $2^\xi$ with
$\xi\geq\lambda$ for $\rho$, and then choose an eligible $\kappa$
above $\rho^+$ and all remaining bounds. Extend $X$ by a presentation
arrow to a level-$\kappa$ object $X_0$, and choose $Y_0\geq X_0$
containing a realization $b$ of the free lift of $p$.

We construct continuous pairs
\[
 X_i\leq Y_i\qquad(i\leq\rho^+)
\]
at level $\kappa$, with a single fixed $b$ in all the upper objects.
The invariant on the marked marginal is
\begin{equation}\label{iso:matrix-marginal}
 \tp(b/X_i)=\tp(b/X_0)^{\mathrm{fr}}_{X_i}.
\end{equation}
For $i<\rho^+$, we also keep increasing continuous sets
\[
 R_i\subseteq X_i,\qquad S_i\subseteq Y_i,\qquad
 R_i\subseteq S_i,\qquad |R_i|,|S_i|\leq\rho,
\]
with $I_p\subseteq R_i$ and $b\subseteq S_i$. They are bookkeeping
sets, not additional structure on the objects.

Here are the two closure requirements at a stage $i$ with
$\cf(i)>\lambda$.
\begin{enumerate}
\item Any specified at-most-$\lambda$ sets of base and top parameters
from $R_i,S_i$ are contained in a pair $A\leq V$ of level-$\lambda$
subobjects of $X_i\leq Y_i$, with $A\subseteq R_i$ and
$V\subseteq S_i$.
\item Every tuple $d$ of length at most $\lambda$ from $S_i$ has a
sequence of at least $\lambda$ copies inside $R_i$ computing its
actual type over $X_i$ on all tests of size less than $\lambda$.
For every subtuple and each such parameter test $Z$, all but at most
$|Z|+\aleph_0$ of the copies give its joint $D$-diagram.
\end{enumerate}

We explain how to arrange the requirements; in particular, they do not
amount to assuming a small saturated submodel. At successor stages,
process all at-most-$\lambda$ lists in the current parameter sets.
There are at most $\rho^\lambda=\rho$ of them. Mixed
L\"owenheim--Skolem supplies the requested level-$\lambda$ pairs.
Lemma~\ref{calc:approximations} supplies copies of each top list in the
current base, at the test sizes below $\lambda$ and with the displayed
exception estimate. Their total number of coordinates is at most
$\rho$. Close the sets under these operations for $\lambda^+$ steps;
the size is still $\rho$. If new elements are added later, repeat these
successor closures. At a limit $i$ of cofinality greater than $\lambda$,
every requested list was contained at an earlier stage. Local character
applied to that list, whose length is at most $\lambda<\cf(i)$, says
that its actual type over $X_i$ is free from some earlier stage.
Approximations placed in a later intervening successor base still
compute it over $X_i$ by the free-extension theorem. Their copies lie
in $R_i$. The same cofinality argument places every small-hull request
in an earlier closure. This proves both requirements at the indicated
limits.

At any such stage, if some list $d$ of length at most $\lambda$ from
$S_i$ makes $\tp(b,d/X_i)$ nonisolated relative to $b$, apply
Proposition~\ref{iso:alternative}. It gives a same-level extension of
$X_i$ in which this joint type has a nonfree extension while the
$b$-marginal remains free. By type-identifying amalgamation, extend
$Y_i$ into a common $Y_{i+1}$, identifying the old $(b,d)$ with that
realization. The new base is the extension just chosen. Transitivity
preserves~\eqref{iso:matrix-marginal}. At other stages use compatible
extensions and the bookkeeping closures. At limits use the working
scheme's unions; Definition~\ref{iso:marginal-continuity} preserves
the marginal invariant. All these unions have size $\kappa$, since
$\rho^+<\kappa$.

This process must reach an indicated stage at which every list in the
second closure requirement is isolated jointly with $b$. Suppose not,
and choose a failing list $d_i$ at every stage with $\cf(i)>\lambda$.
The set of these stages is stationary in the regular cardinal
$\rho^+$. Indeed, in any closed unbounded subset of $\rho^+$ take a
continuous increasing sequence of length $\lambda^+$ and its
supremum; the supremum is still below $\rho^+$, is in the closed
set, and has cofinality $\lambda^+$.

Because $S_i$ is continuous and $|d_i|\leq\lambda<\cf(i)$, there is
$j(i)<i$ with $d_i\subseteq S_{j(i)}$. Fodor's pressing-down argument \cite{Fodor1956} makes
$j(i)$ constant on a stationary subset. For completeness, if every
fiber of a regressive map on a stationary subset of a regular cardinal
were nonstationary, choose a closed unbounded set avoiding each fiber.
Their diagonal intersection is closed unbounded and gives a
contradiction at a point of the stationary set. The diagonal
intersection is unbounded by taking an increasing countable sequence
whose next point is in all the clubs indexed below the previous point;
the needed intersections of fewer than the regular cardinal many clubs
are closed unbounded, by interleaving their choices and taking a
limit. Thus this application uses only set-indexed choices in ZFC.

There are at most $\rho^\lambda=\rho$ possible lists in the resulting
fixed $S_j$. Partitioning the stationary subset into these $\rho$
possibilities retains a stationary piece: otherwise the intersection
of $\rho<\rho^+$ clubs avoiding the pieces is a contradiction.
Consequently one fixed list $d$ witnesses nonfree extensions at
stationarily many stages. In the final union, the tuple $(b,d)$ has
length at most $\lambda<\rho^+$. Local character makes its type
free from some $X_j$. Its restrictions at every later successor are
then free from their preceding base, contradicting the chosen
stationarily many nonfree extensions.

Fix a stage $i$ at which all the required lists are isolated. We next
shrink, and this step must preserve freeness for the tuples on which
isolation will be transferred. Construct increasing pairs
\[
 A_\alpha\leq U_\alpha\qquad(\alpha<\theta^+)
\]
of level-$\lambda$ subobjects inside $R_i,S_i$. Keep $I_p$ in every
base and $b$ in every top. At successor stage $\alpha+1$, take a tuple
enumerating $U_\alpha$. By the second closure requirement it has
$\lambda$ approximation copies in $R_i$. Include every coordinate
of these copies in $A_{\alpha+1}$, and include $A_\alpha,U_\alpha$
in the next pair, using the first closure requirement. The requests
have cardinality $\lambda$. At proper limits take unions; the upper
bound clause into $X_i,Y_i$ keeps them strong there. Finally put
\[
 A=\bigcup_{\alpha<\theta^+}A_\alpha,
 \qquad U=\bigcup_{\alpha<\theta^+}U_\alpha.
\]
Because $\lambda$ is strong limit and $\theta<\lambda$, we have
$\theta^+<\lambda$. Both unions are level-$\lambda$ objects, and
the common upper bound clause gives $A\leq U$.

Let $c$ have length at most $\theta$ in $U$. Regularity of $\theta^+$
puts $(b,c)$ inside some $U_\alpha$. The next base contains $\lambda$
copies of its type over $X_i$, with the exception estimates above.
For any fixed sufficient testing arity below $\lambda$, select a
bounded subsequence longer than that arity's exception bound. These
copies are inside $A$ and compute the actual large-base type. The
majority characterization of free extension now says that
\[
 \tp(b,c/X_i)=\tp(b,c/A)^{\mathrm{fr}}_{X_i}.
\]
Explicitly, move a large-base query into $A$ over the selected copies
and over canonical approximations for the small type. Choose an index
outside the exception sets for the query and its moved copy. The
actual answer therefore equals the small free evaluation. The
testing theorem identifies the whole type. The bounds needed for
this argument are below $\lambda$, since $|b,c|\leq\theta$.

The large type is isolated by the choice of $i$. Its uniform fragment
has size at most $h(|b,c|+\aleph_0)<\lambda$, so
Lemma~\ref{iso:downward-transfer} gives isolation over $A$. Finally,
$I_p$ still determines the marginal on $b$: it did so at $X_0$, the
matrix preserved its free lift, and the same approximation answers
hold on restriction to $A$. Bounded diagram homogeneity gives an
isomorphism $A\simeq X$ carrying these anchors to themselves.
Transport of approximation evaluations carries $\tp(b/A)$ to $p$.
Renaming $U$ along this isomorphism gives the asserted extension of
the original $X$.
\end{proof}

\subsection{All bounded lengths at one level}

The preceding construction isolates tuples through one prescribed length.
All-domination requires this for every bounded tuple in a single object.
We obtain it by descending from a larger level with one common family of
approximate retractions; the argument therefore also covers tuples that
span the countable descent construction.

\begin{theorem}\label{iso:all-lengths}
At sufficiently closed eligible levels, every finite-tuple type
$p=\tp(b/X)$ has an all-dominated realization: there is $Y\geq X$
all-dominated over $X+b$ realizing $p$. More generally this holds for
a bounded marginal satisfying Definition~\ref{iso:marginal-continuity}.
An all-dominated realization embeds over $X,b$ into any actual ambient
containing a realization of that type, with the mixed-level conventions
of Proposition~\ref{iso:prime}.
\end{theorem}

\begin{proof}
Write $\nu$ for the desired level and fix sufficient bounded anchors
$I_p\subseteq X$ for the input marginal. Choose a much larger eligible
$\lambda>\nu$ and lift the type to a level-$\lambda$ base by a
presentation arrow. Apply Proposition~\ref{iso:finite-existence} at
level $\lambda$ with $\theta=\nu$. We obtain a pair $A\leq U$ of
level $\lambda$ such that every tuple of length at most $\nu$ in $U$,
jointly with $b$, has isolation relative to $b$ over $A$. The original
anchors may be retained in $A$ throughout that construction.

Choose $\nu$ distinct indices from the family in
Corollary~\ref{calc:retractions}, and write the corresponding
approximate retractions as
\[
 r_j:U\longrightarrow A\qquad(j<\nu)
\]
such that each $a\in A$ is fixed by all but finitely many of them.
Use mixed L\"owenheim--Skolem to choose level-$\nu$ subobjects
$A'\leq U'$ of $A\leq U$ which contain $I_p,b$ in their indicated
places and satisfy
\begin{equation}\label{iso:retraction-closure}
 r_j[U']\subseteq A'\qquad(j<\nu).
\end{equation}
Here is the closure construction. Start with level-$\nu$ sections
containing the desired entries. From the current top of size $\nu$,
adjoin to the next base all its images under the $\nu$ selected maps;
there are at most $\nu\cdot\nu=\nu$ coordinates. Enlarge the next
top to contain that base and the old top. Iterate countably many
times and take unions. Same-level chain closure and the common upper
bound clause into $A,U$ give the pair and
\eqref{iso:retraction-closure}.

Fix any bounded tuple $d$ in $U'$. The tuple $(b,d)$ is still bounded
at $\nu$. The maps in~\eqref{iso:retraction-closure} place all its
images in $A'$. For a test $Z\subseteq A$, all but at most
$|Z|+\aleph_0$ of the maps fix $Z$ pointwise, and hence for all those
indices
\[
 D(r_j(b,d),Z)=D(b,d,Z).
\]
At a fixed testing arity below $\nu$, choose a bounded subsequence of
indices longer than the exception bound. Lemma~\ref{calc:retraction-descent}
gives
\[
 \tp(b,d/A)=\tp(b,d/A')^{\mathrm{fr}}_A.
\]
The type on the left is isolated, by the length-$\nu$ assertion at
level $\lambda$. Its uniform fragment bound for the actual arity
$|b,d|<\nu$ is below $\nu$. Downward transfer gives isolation of
$\tp(b,d/A')$ relative to $b$.

This argument applies separately to every bounded $d\subseteq U'$;
it makes no assumption that $d$ lies in one stage of the countable
construction. Thus $U'$ is all-dominated over $A'+b$. The retained
anchors and bounded diagram homogeneity identify its marginal with
the original $p$ by an isomorphism $A'\simeq X$. Transport gives the
required object over $X$. Its embedding into an actual ambient is
Proposition~\ref{iso:prime}.
\end{proof}

\subsection{Independent lists and coherent witnesses}

Finite marks already have all-dominated extensions.  To obtain the same
conclusion for a bounded independent list of finite marks, we must first
show that its finite product data determine one whole type and that this
type persists along the actual coupled chains used above.

\begin{definition}\label{iso:independent-definition}
A list $(b_i)_{i<\delta}$ of finite tuples is
\emph{product-independent over $A$} if each finite sublist has the
product of its marginal types over $A$.
A \emph{witnessing chain} is a continuous chain $(M_i)_{i\leq\delta}$
of objects, with $M_0=A$, such that $b_i$ lies in $M_{i+1}$ and
\[
 \tp(b_i/M_i)=\tp(b_i/A)^{\mathrm{fr}}_{M_i}.
\]
It is an \emph{all-dominated-increment chain} if, in addition,
$M_{i+1}$ is all-dominated over $M_i+b_i$. In uses at level $\mu$,
the chain has total size $\mu$, in particular $|\delta|\leq\mu$.
\end{definition}

\begin{lemma}\label{iso:finite-domination}
Suppose $(M_i)_{i\leq\delta}$ is an all-dominated-increment chain over
$A$ for finite tuples $(b_i)_{i<\delta}$. If a finite tuple $x$ in an
actual coupling has product over $A$ with every finite selection of
the accumulated $b_i$, then $\tp(x/M_\delta)$ is free from $A$.
\end{lemma}

\begin{proof}
Induct on the stage. At a successor $i+1$, the finite tuple $(x,b_i)$
has product with every finite selection of earlier marks, by the
finite product rules. The induction hypothesis applied to this tuple
says that its joint type over $M_i$ is free from $A$. Products lift
freely by Proposition~\ref{calc:products}; hence $x$ and $b_i$ still have
product over $M_i$. Proposition~\ref{iso:prime}, applied to the
all-dominated increment $M_{i+1}$, makes $x/M_{i+1}$ free from $M_i$.
Transitivity makes it free from $A$.

At a limit $j$, local character for the finite tuple $x$ gives a
proper stage $i<j$ from which $x/M_j$ is free. The induction
hypothesis and transitivity give freeness from $A$. All types are
computed in the actual coupled ambient, so this argument asserts
continuity of the specified realization, not merely existence of
some independent realization at the limit.
\end{proof}

\begin{lemma}\label{iso:witnessed-uniqueness}
The following hold at sufficiently closed eligible levels.
\begin{enumerate}
\item Every actual product-independent list of finite tuples admits
an all-dominated-increment witnessing chain, retaining its whole
joint configuration in an extension of the given ambient.
\item Two lists supplied with witnessing chains have a unique whole
amalgamation type over $A$, determined by the successive marginal
types over $A$. This assertion also permits bounded, rather than
finite, individual tuples in the witnessing chains. For a list of
total length at most the level, the assertion means that one
amalgam identifies the entire lists; for a bounded total list it is
an equality of the bounded-list types already defined.
\end{enumerate}
\end{lemma}

\begin{proof}
For the first assertion, maintain a common ambient containing the
original list and the already constructed witnessing chain. At stage
$i$, finite domination makes the actual $b_i$ free over $M_i$ from
$A$: its products with finite selections of previous marks are still
the original ones, since arrows preserve types. Theorem~\ref{iso:all-lengths}
supplies an all-dominated realization over $M_i+b_i$. Primeness embeds
it into a compatible actual ambient over this tuple; alternatively,
type-identifying amalgamation supplies such an ambient while retaining
the original list. At limits take the coherent unions. Finite
domination continues to apply, and the chain has the stipulated
total size.

For uniqueness, let $M_i^0,M_i^1$ be the two witnessing chains. Build
common objects $N_i$ and coherent arrows from $M_i^0,M_i^1$ into
$N_i$, agreeing on $A$ and identifying all marks already processed.
At a successor, take the free lift over $N_i$ of the next common
marginal type over $A$, with realization $c$. Its restriction to
either image of $M_i^r$ is the prescribed type of $b_i^r$, by
transitivity. Amalgamate the next stage $M_{i+1}^0$ with this
realization over its old stage, identifying $b_i^0$ with $c$. Then
amalgamate $M_{i+1}^1$ in the same way, identifying its mark with the
same $c$. The first amalgamation preserves the type of $c$ over the
other old-stage image, so the second identification is legitimate.
Retain the old common object throughout. At limits, coherent unions
give the required arrows. The final common ambient identifies every
mark. Only bounded individual marginal types were used in a single
step, which proves the stated version for bounded individual tuples
as well.
\end{proof}

\begin{theorem}\label{iso:independent-lists}
Let $b=(b_i)_{i<\delta}$ be a bounded product-independent list of
finite tuples over $X$ at a sufficiently closed eligible level.
\begin{enumerate}
\item Its whole type is determined by the product data on finite
sublists. Its free lift is product-independent and is the unique
whole type with the lifted finite data.
\item The marked marginal $\tp(b/X)$ satisfies
Definition~\ref{iso:marginal-continuity}.
\item There is an all-dominated object over $X+b$, and it embeds
over these data into an actual ambient containing the given list.
\end{enumerate}
\end{theorem}

\begin{proof}
Lemma~\ref{iso:witnessed-uniqueness} supplies witnesses for any actual
list with the specified finite products and identifies their whole
types. The finite projections of a free lift are the free lifts of
the finite projections, by Proposition~\ref{calc:products}. They are
therefore the required products of the lifted marginal types.
Applying witnessed uniqueness at the receiving base proves the
first assertion, including its mixed-level interpretation.

For continuity, take the actual coupled chain in
Definition~\ref{iso:marginal-continuity}. Each finite sublist of the
fixed mark is a finite tuple. Its actual free evaluations at the
proper stages persist at the union by local character and
transitivity. These finite data are precisely the finite data of
the free lift of the original whole marginal. By the first
assertion they determine its whole type at the union. This proves
the second assertion without applying finite-tuple local character
directly to the possibly long mark.

We may now apply the general forms of
Proposition~\ref{iso:finite-existence} and
Theorem~\ref{iso:all-lengths}. To check the additional arity
conditions, choose the auxiliary $\theta$ at least $|b|$. The
parameter-set scheduling then handles $(b,d)$ at length at most
$\lambda$, and the final stationary contradiction uses length
$\lambda<\rho^+$. At limits of the matrix construction, the
continuity just proved preserves the fixed marked marginal.
All approximation and isolation bounds are evaluated at the actual
bounded joint arities. The all-length descent uses the same
retractions on the concatenations $(b,d)$, which remain bounded at
the desired level. Thus all hypotheses of the two existence
results hold. Primeness gives the assertion about the actual
ambient.
\end{proof}
